\documentclass[10pt,a4paper]{amsart}
\usepackage[margin=19mm]{geometry}
\usepackage{amsmath,amssymb,mathtools}
\usepackage[scr=rsfs,cal=euler]{mathalfa}
\usepackage{xfrac}
\usepackage[unicode,hidelinks,bookmarks=false]{hyperref}
\newcommand{\ie}{i.e.\ }
\newcommand{\cf}{cf.\ }
\newcommand{\resp}{respectively\ }
\newcommand{\Subsection}{\subsection}
\providecommand{\nobreakdash}{\nobreak\hbox{-}}
\setlength{\emergencystretch}{2em}
\multlinegap0pt
\usepackage{bm}
\usepackage{bbm}
\usepackage{dsfont}
\usepackage{xspace}
\usepackage{cancel}
\usepackage{mathtools}
\usepackage{amsthm}
\makeatletter
\@ifundefined{thm}{\newtheorem{thm}{Thm}}{}
\@ifundefined{prop}{\newtheorem{prop}[thm]{Proposition}}{}
\makeatother
\newcommand{\rk}{\mathrm{rk}}
\title[Endomorphisms of free Steiner quasigroups]{Endomorphisms of free Steiner quasigroups}
\author{Silvia Barbina, Enrique Casanovas}
\date{Source: arXiv:2412.06055v3 (CC BY 4.0)}
\begin{document}
\maketitle

\noindent\emph{Source and licence.} Endomorphisms of free Steiner quasigroups. Version arXiv:2412.06055v3 (CC BY 4.0), distributed under the Creative Commons Attribution 4.0 International licence, as recorded in the arXiv licence metadata for this exact version and shown on the abstract page https://arxiv.org/abs/2412.06055v3. Full rights notice: licenses/arxiv-2412.06055-CC-BY-4.0.txt.\par\smallskip

\noindent\textit{Editorial scope.} Counted unit: the Prop whose source label is \texttt{May19\_1} in the article (source0.tex lines 381--387, source label \texttt{May19\_1}), delivered together with the article's own complete original proof of it (source0.tex lines 388--405, one \texttt{proof} environment of 18 lines). No mathematics is written, repaired or re-derived: statement and proof are the article's own text line for line. Required context retained uncounted: none. Every definition, display and label of those lines is retained unchanged. Omitted from this package and not redistributed: 1--380, 406--821. Nothing inside a retained excerpt is omitted or altered: every retained body is delivered exactly as the article's lines are written, so no omission receipt of any kind accompanies this package. Citations: the retained excerpts carry no citation command at all, so no bibliography entry of the article is transcribed and no reference is redistributed. Reference adaptation: none was needed and none was made; every reference inside the delivered unit resolves inside it, so no unresolved reference or citation is produced. Printed slips kept literal: no printed word, symbol, hypothesis or number of the counted statement or of its proof was altered. Layout and class: the article's own class is \texttt{article} with packages including latexsym, amsmath, amssymb, amsthm, enumitem, ulem, tikz, tkz-euclide, geometry, url, comment; the wrapper is the lane's pinned \texttt{amsart} editorial wrapper at 10pt on A4, which restates the theorem-like environments the retained text needs and adds the article's own macro definitions that the retained text needs (\texttt{\textbackslash rk}), with extra wrapper packages (bm, bbm, dsfont, xspace, cancel, mathtools). The delivered numbering is the unit's own. Assets: no retained span contains a figure environment, a graphics-inclusion command, a PDF-page inclusion, a table or a TikZ picture; no image file is copied into this package, the package declares no asset dependency and no image file is redistributed with it. Licence: version arXiv:2412.06055v3 of this article is distributed under the Creative Commons Attribution 4.0 International licence, as recorded in the arXiv licence metadata for this exact version and shown on the abstract page https://arxiv.org/abs/2412.06055v3. A full rights notice is delivered at licenses/arxiv-2412.06055-CC-BY-4.0.txt. Characters: every retained excerpt is pure ASCII, so the delivered engine, XeLaTeX with its default Latin Modern text font, reports no missing character.
	\begin{prop}\label{May19_1}
		Let $M=\langle \bar{a}\rangle$ be a free Steiner quasigroup with standard free construction $(S_i\mid i<\omega)$, let  $\bar{a}= \langle a_i : i \in I \rangle$ be  independent, and let $t= t(\bar{x})$ be a reduced term of rank $k$. Then
		\begin{enumerate}
			\item $t(\bar{a})\in S_k\smallsetminus S_{<k}$
			\item if $r$ is reduced and  $r(\bar{a})= t(\bar{a})$, then $r\sim t$.
		\end{enumerate}
	\end{prop}

	\begin{proof} We prove 1 and 2 jointly by induction on $k$, but with the extra assumption that $\rk(r)=\rk(t)$. This assumption can be eliminated  later since 1 excludes other possibilities.
		
		\underline{Base case $k=0$}.  If $\rk(t)=0$, then $t$ is a variable $x_i$ and  $t(\bar{a})= a_i\in S_0$.  Moreover, if $r$ has rank 0, $r$ is a variable $x_j$  and  $t(\bar{a})= r(\bar{a})$  implies $a_i =a_j$ and therefore  $x_i = x_j$. Hence  $t\sim r$.
		
		\underline{Inductive step $k+1$}.  Let $t= t_1\cdot t_2$, with $k_1= \rk(t_1)\leq \rk(t_2) = k$ and let  $b= t(\bar{a})$.  By inductive hypothesis, $b_1= t_1(\bar{a})\in S_{k_1}\smallsetminus S_{<k_1}$ and $b_2= t_2(\bar{a})\in S_k\smallsetminus S_{<k}$. Then $b= b_1\cdot b_2$ and there are different cases.
		
		\emph{Case 1}. Assume $k_1 = k$.  If  $b_1 = b_2$, then $b= b_1= b_2$ and by inductive hypothesis $t_1\sim t_2$. But then $t$ is not reduced. Hence, 
		$b_1\neq b_2$. By construction, $b\in S_{k+1}\smallsetminus S_k$.  Moreover, if $r$ is reduced of rank $k+1$ with $r(\bar{a})= b$, then $r= r_1\cdot r_2$, with $r_1,r_2$ reduced and, without loss of generality, $\rk(r_1)\leq \rk(r_2)= k$. By inductive hypothesis $b_1^\prime = r_1(\bar{a})\in S_{\rk(r_1)}\smallsetminus S_{<\rk(r_1)}$  and $b_2^\prime = r_2(\bar{a})\in S_k\smallsetminus S_{<k}$. Then $b= b_1^\prime\cdot b_2^\prime$ and by construction $\{b_1,b_2\}= \{b_1^\prime,b_2^\prime\}$. If $b_1= b_1^\prime$ and $b_2 = b_2^\prime$, then $\rk(r_1)= k_1$ and by inductive hypothesis $t_1\sim r_1$ and $t_2\sim r_2$. Therefore in this case $t\sim r$. If $b_1= b_2^\prime$ and $b_2= b_1^\prime$, then $k=k_1= \rk(r_1)$ and, by inductive hypothesis, $t_1\sim r_2$ and $t_2\sim r_1$. Again we conclude that $t\sim r$.
		
		\emph{Case 2}. Assume $k_1 <k$ and $b\not \in S_{<k}$.  Note that $b\neq b_2$ and $b\not \in S_k$. By construction, $b= b_1\cdot b_2 \in S_{k+1}$.  Now let $r$ be reduced, of rank $k+1$ and such that $r(\bar{a})= t(\bar{a})=b$. Exactly as in the previous case we conclude that $r\sim t$.
		
		\emph{Case 3}. The last case is $k_1 < k$ and  $b\in S_j\smallsetminus S_{<j}$  for some $j<k$.  We show this is impossible.   We know that there is some reduced term $s$ of rank $j$ such that $s(\bar{a}) = b$. Now we claim that $s\cdot t_1$ is reduced. Otherwise, there are two possible cases:
		\begin{enumerate}
			\item[(i)] $s\sim t_1$.  But then $b= b_1$ and hence $b=b_1 = b_2$, which is not the case.
			\item[(ii)]  $t_1\sim s\cdot r$ or $s\sim t_1\cdot r$  for some (reduced)  term $r$.  Then $r(\bar{a})= t_1(\bar{a})\cdot s(\bar{a})= b_1\cdot b= b_2$  and by inductive hypothesis $\rk(r)=k$, contradicting $\rk(t_1)=k_1<k$ and $\rk(s)= j<k$.
		\end{enumerate}
		Therefore, we conclude that $s\cdot t_1$ is reduced. Then, since $\rk(s\cdot t_1)\leq k$ and $s\cdot  t_1 (\bar{a}) = b\cdot b_1 = b_2$, we can apply the inductive hypothesis to get $\rk(s\cdot t_1)= k$  and  $s\cdot t_1\sim t_2$. Then $t\sim t_1\cdot t_2\sim t_1\cdot (s\cdot t_1)$, which says that $t$ is not reduced.
	\end{proof}

\end{document}
